% P03.1 curvewise scalar-defect localization proof, live integration.

\begin{proof}[\Cref{thm:p031-curvewise-scalar-defect-localization}]
For every $i$, scalar-extension stability and analytic-pullback stability
\cite[Proposition~4.2(5),(6)]{GR25Classical} make
$u_i\circ\pi\circ g$ classically plurisubharmonic, equivalently
subharmonic, on $C$.  The pulled family remains locally upper bounded, and
its pointwise supremum is $s\circ\pi\circ g$.  Therefore
\Cref{lem:p031-rigsmooth-curve-envelope} shows that $p_g$ is subharmonic.

Applying \Cref{thm:p031-affine-ground-field-perron} over $K'$ to the
family $(u_i\circ\pi)_{i\in I}$ shows that
$v=(s\circ\pi)^*$ is $K'$-curvewise plurisubharmonic.  Hence $h_g=v\circ g$
is subharmonic by \Cref{def:p031-ground-field-curvewise-psh}.

We next prove the two inequalities in
\eqref{eq:p031-curvewise-scalar-defect-chain}.  If $a:X\to
\mathbf R\cup\{-\infty\}$ is locally bounded above and $f:Y\to X$ is
continuous, then
\begin{equation}
\label{eq:p031-v61-continuous-pullback-usc-inequality}
                         (a\circ f)^*\leq a^*\circ f.
\end{equation}
Indeed, for $y\in Y$ and every neighborhood $U$ of $f(y)$,
\[
        \sup_{f^{-1}(U)}a\circ f\leq\sup_U a.
\]
Taking the neighborhood-filter infima that define upper-semicontinuous
regularization proves
\eqref{eq:p031-v61-continuous-pullback-usc-inequality}.  Apply it to
$a=s\circ\pi$ and $f=g$ to obtain $p_g\leq h_g$.

The canonical-retraction normal form
\eqref{eq:p031-canonical-retraction-regularization}, applied to $a=s$,
gives
\[
                         v\leq v\circ r,
                 \qquad u\circ\pi=v\circ r.
\]
Composition with $g$ yields $h_g\leq q_g$ and the displayed formula for
$q_g$.  The chain $p_g\leq h_g\leq q_g$ immediately implies
\eqref{eq:p031-curvewise-total-defect-splitting}.

If $c\notin\overline{\mathcal B_g}$, there is an open neighborhood $U$ of
$c$ disjoint from $\mathcal B_g$.  Since $h_g\leq q_g$, the definition of
$\mathcal B_g$ gives $q_g|_U=h_g|_U$.  Thus $q_g$ is subharmonic on $U$.
The sheaf property for subharmonic functions
\cite[Proposition~3.5(1)]{GR25Classical} proves subharmonicity on the union
$C\setminus\overline{\mathcal B_g}$.  Taking the contrapositive gives
\eqref{eq:p031-curvewise-failure-contained-in-scalar-defect}.
\end{proof}

\begin{proof}[\Cref{cor:p031-scalar-defect-curve-witness}]
By the all-extensions smooth strict curve test
\cite[Definition~4.1 and \S4.1.3]{GR25Classical}, failure of classical
plurisubharmonicity supplies an analytic extension $K'/K$, a smooth
connected strict $K'$-analytic curve $C_0$, and an analytic morphism
$g_0:C_0\to V_{K'}$ for which $u\circ\pi_{K'/K}\circ g_0$ is not
subharmonic.  Sheafiness then gives a point $c$ at which this function is
not subharmonic on any neighborhood.  Analytic curves are separated under
the standing convention of \cite[\S3.1]{GR25Classical}.  We may therefore
apply \eqref{eq:p031-curvewise-failure-contained-in-scalar-defect} directly
with $C=C_0$ and $g=g_0$.  Since
\[
 h_g=(s\circ\pi_{K'/K})^*\circ g,
 \qquad
 q_g=u\circ\pi_{K'/K}\circ g,
\]
membership of $c$ in $\overline{\mathcal B_g}$ is exactly the displayed
accumulation condition.
\end{proof}
